\chapter{Введение}
\label{ch:intro}

\section{Введение}

На прошлом занятии мы формировали семплы, используя BPSK/QPSK модуляцию. Мы использовали формирующий фильтр, имеющий
импульсную характеристику прямоугольной формы. Прямоугольная импульсная характеристика не используется в реальном оборудовании, потому 
что прямоугольный сигнал будет занимать слишком много спектра. На этом занятии зададим для формирующего фильтра импульсную характеристику
с формой приподнятого косинуса. Такая форма может использоваться в реальном оборудовании.

\section{Цель}

Осуществить формирование и передачу QPSK символов, используя формирующий фильтр, импульсная
характеристика которого имеет форму приподнятого косинуса.


\endinput